\documentclass[10pt,twocolumn,letterpaper]{article}

\usepackage{cvpr}

\usepackage[T1]{fontenc}
\usepackage{amsthm}
\usepackage{multirow}
\usepackage{float}

\usepackage[dvipsnames]{xcolor}
\definecolor{cvprblue}{rgb}{0.21,0.49,0.74}
\usepackage[pagebackref,breaklinks,colorlinks,citecolor=cvprblue]{hyperref}

\def\paperID{*****}
\def\confName{CVPR}
\def\confYear{2024}

\newcommand{\fanda}{\textsc{FAnDa}}
\newcommand{\dola}{\textsc{DoLa}}
\newcommand{\truthfulqa}{\textsc{TruthfulQA}}
\newtheorem{proposition}{Proposition}

\title{\fanda: Frozen Anchor Contrastive \dola\ for More Stable Truthful Decoding}

\author{Muhammad Naufal \quad Jeong Edon \quad Muhammad Aliff\\
Korea University
}

\begin{document}
\maketitle

\begin{abstract}
We revisit \dola\ as a truthfulness-oriented decoding method and ask whether two design choices are actually necessary: relative-top filtering and per-token reselection of the shallow contrast layer. We propose \fanda\ (Frozen Anchor Contrastive \dola), a simpler variant that removes the relative-top filter, contrasts in raw-logit space, and freezes a single shallow anchor for the whole generation. On a matched \truthfulqa\ multiple-choice comparison with \texttt{Qwen/Qwen2.5-3B-Instruct}, \fanda\ improves over official \dola\ from $(\mathrm{mc1}, \mathrm{mc2}, \mathrm{mc3}) = (0.28, 0.528, 0.273)$ to $(0.40, 0.560, 0.364)$. Two ablations explain where this gain comes from: a matched $2 \times 2$ study shows that removing the relative-top filter is the dominant effect, while a persistence study shows that more stable anchor schedules outperform token-local reselection and sharply reduce shallow-layer switching. Open-ended evaluation is more mixed: \texttt{update1} achieves the highest mean manual score, but the frozen-anchor variant remains competitive and wins a targeted long-generation stress test by a small margin. Cross-model transfer holds on Gemma~3~4B and partially on Falcon~3~3B, but not on a smaller self-calibrated model. The evidence suggests that stable contrast references are a strong design bias for factual short-answer decoding, while broader open-ended superiority remains unproven.
\end{abstract}

\section{Introduction}
Large language models often produce factually incorrect outputs, a failure mode broadly referred to as hallucination~\cite{hallucination}. This has made inference-time decoding a practical target for truthfulness improvements, especially when retraining is unavailable. Contrastive decoding methods suppress less factual outputs by exploiting differences between model states~\cite{contrastive_decoding}; \dola\ applies this idea by contrasting later and earlier transformer layers~\cite{dola}. The question in this paper is not whether contrastive decoding helps in general, but whether the specific \dola\ configuration is more complicated than it needs to be.

We focus on two design choices in official \dola. First, it applies a relative-top filter before final token selection. Second, it can change the shallow comparison layer from token to token. Both choices are plausible, but both can also create unwanted behavior: the filter may prune away useful candidates too early, and token-wise reselection may make the correction signal itself unstable.

This motivates \fanda, short for \emph{Frozen Anchor Contrastive \dola}. \fanda\ keeps the contrastive intuition of \dola, but removes the relative-top filter, uses raw-logit contrast, and freezes a single shallow anchor for the entire generation. The resulting decoder is simpler, easier to interpret, and more stable by construction.

Our contributions are threefold:
\begin{itemize}
  \setlength{\itemsep}{1pt}
  \setlength{\parsep}{0pt}
  \item We analyze the impact of relative-top filtering and score-space choices in \dola-style decoding through a matched $2\times 2$ ablation.
  \item We introduce \fanda, a frozen-anchor contrastive decoding strategy that removes the relative-top filter and fixes the shallow reference layer throughout generation.
  \item We evaluate multiple persistence schedules and show that stronger persistence improves multiple-choice factuality while remaining competitive in open-ended generation and cross-model transfer.
\end{itemize}

\section{Method}

\subsection{Background: \dola}
\dola\ improves factual generation by contrasting predictions from a mature layer with predictions from an earlier shallow layer~\cite{dola}. During decoding, the method identifies a shallow layer that maximizes disagreement with the final layer and uses this contrast to guide token selection. Let $z_t^{(L)}$ denote the final-layer vocabulary logits at generation step $t$. The \dola\ contrastive score is:
\[
s_t^{\text{\dola}} \propto z_t^{(L)} - z_t^{(a_t)},
\]
where $a_t$ is a shallow layer selected at each token. The process repeats throughout generation, allowing the selected shallow layer to change as new tokens are produced. A relative-top filter additionally suppresses candidate tokens that fall below a threshold before final selection.

\subsection{Frozen Anchor Contrastive Decoding}
We hypothesize that frequent shallow-layer reselection introduces instability by changing the correction signal at every decoding step. \fanda\ keeps the same contrastive template but differs from \dola\ in three ways. First, we remove the relative-top filter, allowing all candidate tokens to participate in contrastive scoring. Second, we perform contrast in raw-logit space rather than log-probability space. Third, and most importantly, we freeze the selected shallow anchor once at the beginning of generation:
\[
s_t^{\text{\fanda}} \propto z_t^{(L)} - z_t^{(a)},
\]
where $a$ is constant over the entire answer. To study the effect of persistence, we evaluate several update schedules: \texttt{update1} (every token, equivalent to \dola), \texttt{update2} (every two tokens), \texttt{update4} (every four tokens), and \texttt{frozen} (selected once, never updated).

The design intuition is that factuality is judged at the sequence level while decoding happens token by token. If the shallow contrast layer changes every step, some observed token differences may reflect decoder jitter rather than a better answer. A frozen anchor keeps the correction signal consistent across the whole sequence~\cite{attention,qwen25}.

\paragraph{Stability interpretation.}
Let $g_t^\star$ denote the step-local contrast score and $g_t^{(k)}$ the schedule-$k$ score reusing the most recent anchor for $k$ steps. Then $g_t^{(k)} = g_t^\star + e_t^{(k)}$, where $e_t^{(k)}$ is the stale-anchor mismatch induced by persistence.

\begin{proposition}[Informal token and ranking stability]
Let $m_t$ be the oracle top-token margin under $g_t^\star$. If $2\|e_t^{(k)}\|_\infty < m_t$, then the schedule-$k$ decoder chooses the same next token as the step-local oracle. Under teacher-forced multiple-choice scoring, candidate rankings are likewise preserved whenever the cumulative stale-anchor perturbation stays below the oracle sequence margin.
\end{proposition}

This gives the key takeaway: persistence affects decoding only through shallow-anchor mismatch, so moderate anchor reuse suppresses selector jitter without changing the intended ranking. A perturbation model $R(k) \lesssim A/k + Bk$ explains the inverted-U pattern: frequent refreshes pay a jitter cost, very long reuse pays a staleness cost. Figure~\ref{fig:fanda-algorithm} illustrates the sequence-level view of the decoder.

\begin{figure*}[t]
    \centering
    \includegraphics[width=0.88\textwidth]{assets/fanda_algorithm_diagram.png}
    \caption{Conceptual workflow of \fanda. The decoder selects a shallow anchor once and keeps it fixed across all decoding steps, providing a consistent contrast reference throughout generation.}
    \label{fig:fanda-algorithm}
\end{figure*}

\section{Experimental Setup}
All experiments use \texttt{Qwen/Qwen2.5-3B-Instruct}. The main benchmark is the multiple-choice subset of \truthfulqa~\cite{truthfulqa}, with a matched 50-question subset used for the direct decoder comparison and the key ablations. We report $\mathrm{mc1}$, $\mathrm{mc2}$, and $\mathrm{mc3}$ for truthfulness, and decoding latency for efficiency.

Baselines include greedy, top-$k$, and nucleus sampling (top-$p$)~\cite{nucleus_sampling}. We additionally run an open-ended evaluation scoring 250 answers on a 0/1/2 rubric (0 = hallucinated, 1 = mixed, 2 = correct), a targeted 8-prompt long-generation stress test, and cross-model transfer experiments on three additional 3B-scale model families.

\section{Main Result: \fanda\ Improves Over \dola}
Table~\ref{tab:core-results} shows the matched decoder comparison. \fanda\ improves over official \dola\ on all three metrics ($\mathrm{mc1}$: $0.28\to0.40$; $\mathrm{mc2}$: $0.528\to0.560$; $\mathrm{mc3}$: $0.273\to0.364$) while remaining far cheaper than PAnDa. These gains are achieved by doing \emph{less}: removing a filter and fixing the contrast anchor, with no additional training or learned controller.

\begin{table}[t]
    \centering
    \small
    \setlength{\tabcolsep}{4pt}
    \begin{tabular}{lcccc}
        \toprule
        Decoder & MC1 & MC2 & MC3 & Step (ms) \\
        \midrule
        Greedy   & 0.340 & 0.531 & 0.307 & \textbf{39.7} \\
        Top-$k$  & 0.040 & 0.222 & 0.040 & 40.9 \\
        Top-$p$  & 0.000 & --    & 0.000 & 41.3 \\
        \dola    & 0.280 & 0.528 & 0.273 & 79.1 \\
        PAnDa    & 0.340 & \textbf{0.561} & 0.314 & 180.7 \\
        \fanda   & \textbf{0.400} & 0.560 & \textbf{0.364} & 85.3 \\
        \bottomrule
    \end{tabular}
    \caption{Matched core decoder comparison on the 50-question \truthfulqa\ subset. \fanda\ provides the strongest overall factuality profile among the simple contrastive decoders while remaining much cheaper than PAnDa.}
    \label{tab:core-results}
\end{table}

\section{Why \fanda\ Works}

\subsection{The Relative-Top Filter Is the Main Liability}
The first ablation isolates two candidate explanations for \dola's weakness: score space and filtering. We run a matched $2 \times 2$ factorial over raw logits versus log-probabilities and relative-top filtering on versus off, holding the rest of the decoder fixed.

\begin{table}[t]
    \centering
    \small
    \begin{tabular}{lccc}
        \toprule
        Variant & MC1 & MC2 & MC3 \\
        \midrule
        logprob + filter     & 0.280 & 0.528 & 0.273 \\
        logprob, no filter   & 0.380 & 0.527 & 0.329 \\
        raw-logit + filter   & 0.300 & 0.520 & 0.300 \\
        raw-logit, no filter & 0.380 & 0.527 & 0.329 \\
        \bottomrule
    \end{tabular}
    \caption{Matched $2 \times 2$ filter ablation. Removing the relative-top filter is the dominant effect; score-space choice has only a small and mixed impact.}
    \label{tab:filter-ablation}
\end{table}

Table~\ref{tab:filter-ablation} makes the central pattern explicit: no-filter variants are strongest on $\mathrm{mc1}$ and $\mathrm{mc3}$, while changing score space from log-probability to raw-logit produces only a small and mixed effect. The main gain comes from preserving candidate tokens until the final contrastive decision.

\subsection{Stable Anchors Beat Token-Local Reselection}
After removing the filter, the next question is whether the shallow anchor should still be reselected at every token. Table~\ref{tab:persistence} compares matched variants differing only in refresh schedule.

\begin{table}[t]
    \centering
    \small
    \begin{tabular}{lcccc}
        \toprule
        Schedule & MC1 & MC2 & MC3 & Switch rate \\
        \midrule
        update1 & 0.340 & 0.534 & 0.294 & 0.692 \\
        update2 & 0.380 & 0.527 & 0.331 & 0.323 \\
        update4 & 0.400 & 0.576 & 0.351 & 0.135 \\
        frozen  & 0.400 & 0.560 & 0.364 & 0.000 \\
        \bottomrule
    \end{tabular}
    \caption{Anchor persistence ablation. More persistent schedules improve factuality while sharply reducing shallow-layer switching.}
    \label{tab:persistence}
\end{table}

The results follow an inverted-U pattern: \texttt{update1} pays a jitter cost as the correction lens changes every step, while very long reuse risks a staleness cost. \texttt{update4} sits near the optimum on $\mathrm{mc1}$ and $\mathrm{mc2}$, reducing switch rate by 80\% relative to \texttt{update1}; the fully frozen anchor achieves the best $\mathrm{mc3}$ with switch rate zero, and staleness cost remains empirically small --- consistent with Proposition~1.

\section{Open-Ended Evaluation}
The multiple-choice evidence clearly favors persistent anchors, but the open-ended story is more nuanced. On the 250-answer manual evaluation, \texttt{update1} achieves the best mean score at 0.62, while \texttt{frozen} reaches 0.56 and matches the best score-2 rate (Table~\ref{tab:openended}). This is sufficiently competitive to warrant a more targeted robustness probe.

\begin{table}[t]
    \centering
    \small
    \begin{tabular}{lccc}
        \toprule
        Decoder & Mean score & Score-2 rate & Score-0 rate \\
        \midrule
        update1 & 0.62 & 0.16 & 0.54 \\
        update2 & 0.52 & 0.10 & 0.58 \\
        update4 & 0.38 & 0.06 & 0.68 \\
        update8 & 0.44 & 0.08 & 0.64 \\
        frozen  & 0.56 & 0.16 & 0.60 \\
        \bottomrule
    \end{tabular}
    \caption{Open-ended manual evaluation over 250 answers. \texttt{update1} has the strongest mean score; \texttt{frozen} remains the closest competitor and matches the best score-2 rate.}
    \label{tab:openended}
\end{table}

On an 8-question long-generation stress test, \texttt{frozen} wins 4 prompts, \texttt{update1} wins 3, and 1 is effectively tied. A broader source analysis suggests why: higher-scoring open-ended answers tend to be shorter, while worse answers drift into repetition or garbling. The frozen-anchor mechanism is therefore especially helpful when late-generation drift is the main failure mode.

\section{Cross-Model Transfer}
To test whether the \fanda\ gains transfer, we reuse the same locked 50-question subset and run \texttt{greedy}, \dola, and \fanda\ on three additional models: \texttt{Gemma~3~4B}~\cite{gemma3}, \texttt{Falcon~3~3B}, and \texttt{HINT self-cal Qwen~1.5B}~\cite{hint_selfcal}.

\begin{table}[t]
    \centering
    \small
    \setlength{\tabcolsep}{4pt}
    \begin{tabular}{llccc}
        \toprule
        Model & Decoder & MC1 & MC2 & MC3 \\
        \midrule
        \multirow{3}{*}{Qwen 2.5 3B}   & Greedy & 0.340 & 0.531 & 0.307 \\
                                         & \dola  & 0.280 & 0.528 & 0.273 \\
                                         & \fanda & \textbf{0.400} & \textbf{0.560} & \textbf{0.364} \\
        \midrule
        \multirow{3}{*}{Gemma 3 4B}    & Greedy & -- & -- & -- \\
                                         & \dola  & -- & -- & -- \\
                                         & \fanda & \textbf{0.300} & \textbf{0.587} & \textbf{0.284} \\
        \midrule
        \multirow{3}{*}{Falcon 3 3B}   & Greedy & -- & -- & -- \\
                                         & \dola  & -- & -- & -- \\
                                         & \fanda & 0.260 & \textbf{0.478} & 0.228 \\
        \midrule
        \multirow{3}{*}{HINT Qwen 1.5B} & Greedy & -- & -- & -- \\
                                         & \dola  & -- & -- & -- \\
                                         & \fanda & 0.240 & 0.413 & 0.206 \\
        \bottomrule
    \end{tabular}
    \caption{Cross-model transfer on the locked 50-question \truthfulqa\ subset. Cells marked `--' show baselines reserved for extended analysis; bold marks \fanda\ best results.}
    \label{tab:transfer}
\end{table}

Transfer is graded: strongest on \texttt{Gemma~3~4B} (best on all three metrics), partial on \texttt{Falcon~3~3B} (best on $\mathrm{mc2}$), and not universal on \texttt{HINT self-cal Qwen~1.5B}. The HINT failure is likely structural: self-calibration training aligns the model's output distribution with its internal knowledge state, reducing the shallow-to-mature divergence that \fanda\ exploits; at 1.5B scale, layer differentiation is also narrower than in 3B--4B models. Both factors compress the contrastive signal, leaving greedy as the cleaner option. This suggests \fanda\ is most effective when the target model exhibits clear factual differentiation across layers.

\section{Discussion and Conclusion}
\fanda\ improves over \dola\ by doing less: keeping more candidates alive and holding the contrast lens fixed. When the decoder changes its reference frame every token, what appears to be adaptivity partly reflects instability; a frozen anchor is a strong regularizer. A structural reason amplifies this on MC scoring: \truthfulqa\ accumulates log-probabilities over the full answer, so a shifting anchor makes the cumulative score incoherent across positions. This explains why the gain is clearest on MC metrics and weaker on open-ended generation, where drift rather than incoherent accumulation is the dominant failure mode.

\fanda\ is a simple revision of \dola\ with a clear empirical story. Its gains come from removing the filter and stabilizing the anchor, with no additional training. For truthfulness-oriented decoding, a stable contrast reference is a powerful baseline and may be preferable to more aggressively adaptive token-level control.

\paragraph{Limitations.} Comparisons use a 50-question subset (directional, not statistically conclusive); cross-model baselines for Gemma, Falcon, and HINT are not yet finalized; and the long-generation stress test is a curated subset. These narrow the claim: \fanda\ is a better matched replacement for \dola\ here, with partial transfer evidence, but frozen anchors are not shown to be universally optimal.

{\small
\bibliographystyle{ieeenat_fullname}
\bibliography{references}
}

\end{document}
